%!TEX program = xelatex
\documentclass[11pt,a4paper]{article}
\usepackage[margin=2.2cm]{geometry}
\usepackage{amsmath,amssymb,amsthm,mathtools}
\usepackage{bm}
\usepackage{enumitem}
\usepackage{booktabs}
\usepackage{xcolor}
\usepackage{hyperref}
\usepackage{fancyhdr}
\usepackage{xeCJK}

\hypersetup{
    colorlinks=true,
    linkcolor=blue!60!black,
    citecolor=blue!60!black,
    urlcolor=blue!60!black
}

\pagestyle{fancy}
\fancyhf{}
\rhead{\small Qiuzhen QE Computational \& Applied Math --- Fall 2024}
\lhead{\small Official Qualifying Exam}
\rfoot{\small Page \thepage}

\DeclareMathOperator*{\argmax}{arg\,max}
\DeclareMathOperator*{\argmin}{arg\,min}

\setlist[itemize]{topsep=3pt,itemsep=2pt,parsep=0pt}
\setlist[enumerate]{topsep=3pt,itemsep=2pt,parsep=0pt}

\newcommand{\examstatement}[1]{%
  \par\addvspace{0.85em}%
  \noindent{\bfseries #1}\par\nobreak\addvspace{0.28em}%
}

% Shared amsthm note/remark/theorem styles
% Shared math extras for qe-review.
% Requires already-loaded: amsmath, amssymb.
%
% Classic pitfall: AMS blackboard-bold fonts have no digit glyphs, so
% \mathbb{1} renders as overlapping garbage. Map the digit 1 to dsfont's
% \mathds{1}; leave \mathbb{R}, \mathbb{P}, etc. on AMS.
%
% Usage (path relative to the main .tex being compiled):
%   notes:  % Shared math extras for qe-review.
% Requires already-loaded: amsmath, amssymb.
%
% Classic pitfall: AMS blackboard-bold fonts have no digit glyphs, so
% \mathbb{1} renders as overlapping garbage. Map the digit 1 to dsfont's
% \mathds{1}; leave \mathbb{R}, \mathbb{P}, etc. on AMS.
%
% Usage (path relative to the main .tex being compiled):
%   notes:  % Shared math extras for qe-review.
% Requires already-loaded: amsmath, amssymb.
%
% Classic pitfall: AMS blackboard-bold fonts have no digit glyphs, so
% \mathbb{1} renders as overlapping garbage. Map the digit 1 to dsfont's
% \mathds{1}; leave \mathbb{R}, \mathbb{P}, etc. on AMS.
%
% Usage (path relative to the main .tex being compiled):
%   notes:  \input{../../common/qe-math-extras.tex}
%   exams:  \input{../../../common/qe-math-extras.tex}

\usepackage{dsfont}
\usepackage{pdftexcmds}

\makeatletter
\let\qe@amsmmathbb\mathbb
\renewcommand{\mathbb}[1]{%
  \ifnum\pdf@strcmp{\unexpanded{#1}}{1}=\z@
    \mathds{1}%
  \else
    \qe@amsmmathbb{#1}%
  \fi
}
\makeatother

% Preferred indicator macros (either style is safe after the remap above).
\newcommand{\bbone}{\mathds{1}}
\newcommand{\ind}{\mathbf{1}}

%   exams:  % Shared math extras for qe-review.
% Requires already-loaded: amsmath, amssymb.
%
% Classic pitfall: AMS blackboard-bold fonts have no digit glyphs, so
% \mathbb{1} renders as overlapping garbage. Map the digit 1 to dsfont's
% \mathds{1}; leave \mathbb{R}, \mathbb{P}, etc. on AMS.
%
% Usage (path relative to the main .tex being compiled):
%   notes:  \input{../../common/qe-math-extras.tex}
%   exams:  \input{../../../common/qe-math-extras.tex}

\usepackage{dsfont}
\usepackage{pdftexcmds}

\makeatletter
\let\qe@amsmmathbb\mathbb
\renewcommand{\mathbb}[1]{%
  \ifnum\pdf@strcmp{\unexpanded{#1}}{1}=\z@
    \mathds{1}%
  \else
    \qe@amsmmathbb{#1}%
  \fi
}
\makeatother

% Preferred indicator macros (either style is safe after the remap above).
\newcommand{\bbone}{\mathds{1}}
\newcommand{\ind}{\mathbf{1}}


\usepackage{dsfont}
\usepackage{pdftexcmds}

\makeatletter
\let\qe@amsmmathbb\mathbb
\renewcommand{\mathbb}[1]{%
  \ifnum\pdf@strcmp{\unexpanded{#1}}{1}=\z@
    \mathds{1}%
  \else
    \qe@amsmmathbb{#1}%
  \fi
}
\makeatother

% Preferred indicator macros (either style is safe after the remap above).
\newcommand{\bbone}{\mathds{1}}
\newcommand{\ind}{\mathbf{1}}

%   exams:  % Shared math extras for qe-review.
% Requires already-loaded: amsmath, amssymb.
%
% Classic pitfall: AMS blackboard-bold fonts have no digit glyphs, so
% \mathbb{1} renders as overlapping garbage. Map the digit 1 to dsfont's
% \mathds{1}; leave \mathbb{R}, \mathbb{P}, etc. on AMS.
%
% Usage (path relative to the main .tex being compiled):
%   notes:  % Shared math extras for qe-review.
% Requires already-loaded: amsmath, amssymb.
%
% Classic pitfall: AMS blackboard-bold fonts have no digit glyphs, so
% \mathbb{1} renders as overlapping garbage. Map the digit 1 to dsfont's
% \mathds{1}; leave \mathbb{R}, \mathbb{P}, etc. on AMS.
%
% Usage (path relative to the main .tex being compiled):
%   notes:  \input{../../common/qe-math-extras.tex}
%   exams:  \input{../../../common/qe-math-extras.tex}

\usepackage{dsfont}
\usepackage{pdftexcmds}

\makeatletter
\let\qe@amsmmathbb\mathbb
\renewcommand{\mathbb}[1]{%
  \ifnum\pdf@strcmp{\unexpanded{#1}}{1}=\z@
    \mathds{1}%
  \else
    \qe@amsmmathbb{#1}%
  \fi
}
\makeatother

% Preferred indicator macros (either style is safe after the remap above).
\newcommand{\bbone}{\mathds{1}}
\newcommand{\ind}{\mathbf{1}}

%   exams:  % Shared math extras for qe-review.
% Requires already-loaded: amsmath, amssymb.
%
% Classic pitfall: AMS blackboard-bold fonts have no digit glyphs, so
% \mathbb{1} renders as overlapping garbage. Map the digit 1 to dsfont's
% \mathds{1}; leave \mathbb{R}, \mathbb{P}, etc. on AMS.
%
% Usage (path relative to the main .tex being compiled):
%   notes:  \input{../../common/qe-math-extras.tex}
%   exams:  \input{../../../common/qe-math-extras.tex}

\usepackage{dsfont}
\usepackage{pdftexcmds}

\makeatletter
\let\qe@amsmmathbb\mathbb
\renewcommand{\mathbb}[1]{%
  \ifnum\pdf@strcmp{\unexpanded{#1}}{1}=\z@
    \mathds{1}%
  \else
    \qe@amsmmathbb{#1}%
  \fi
}
\makeatother

% Preferred indicator macros (either style is safe after the remap above).
\newcommand{\bbone}{\mathds{1}}
\newcommand{\ind}{\mathbf{1}}


\usepackage{dsfont}
\usepackage{pdftexcmds}

\makeatletter
\let\qe@amsmmathbb\mathbb
\renewcommand{\mathbb}[1]{%
  \ifnum\pdf@strcmp{\unexpanded{#1}}{1}=\z@
    \mathds{1}%
  \else
    \qe@amsmmathbb{#1}%
  \fi
}
\makeatother

% Preferred indicator macros (either style is safe after the remap above).
\newcommand{\bbone}{\mathds{1}}
\newcommand{\ind}{\mathbf{1}}


\usepackage{dsfont}
\usepackage{pdftexcmds}

\makeatletter
\let\qe@amsmmathbb\mathbb
\renewcommand{\mathbb}[1]{%
  \ifnum\pdf@strcmp{\unexpanded{#1}}{1}=\z@
    \mathds{1}%
  \else
    \qe@amsmmathbb{#1}%
  \fi
}
\makeatother

% Preferred indicator macros (either style is safe after the remap above).
\newcommand{\bbone}{\mathds{1}}
\newcommand{\ind}{\mathbf{1}}

% Shared amsthm environments for qe-review (minimal).
% Requires already-loaded: amsthm.
% Safe to \input after \usepackage{amsmath,amssymb,amsthm,...}.
%
% Shared numbering uses aliascnt so cleveref keys off the environment
% name (Lemma, Proposition, Exercise, ...) rather than the parent
% counter (Theorem, Definition, Remark). Without aliascnt, \cref{lem:...}
% prints ``Theorem 9.13'' instead of ``Lemma 9.13''.

\usepackage{aliascnt}

\theoremstyle{plain}
\newtheorem{theorem}{Theorem}[section]

\newaliascnt{lemma}{theorem}
\newtheorem{lemma}[lemma]{Lemma}
\aliascntresetthe{lemma}

\newaliascnt{proposition}{theorem}
\newtheorem{proposition}[proposition]{Proposition}
\aliascntresetthe{proposition}

\newaliascnt{corollary}{theorem}
\newtheorem{corollary}[corollary]{Corollary}
\aliascntresetthe{corollary}

\newaliascnt{claim}{theorem}
\newtheorem{claim}[claim]{Claim}
\aliascntresetthe{claim}

\newaliascnt{fact}{theorem}
\newtheorem{fact}[fact]{Fact}
\aliascntresetthe{fact}

\theoremstyle{definition}
\newtheorem{definition}{Definition}[section]

\newaliascnt{exercise}{definition}
\newtheorem{exercise}[exercise]{Exercise}
\aliascntresetthe{exercise}

\newaliascnt{assumption}{definition}
\newtheorem{assumption}[assumption]{Assumption}
\aliascntresetthe{assumption}

\newaliascnt{notation}{definition}
\newtheorem{notation}[notation]{Notation}
\aliascntresetthe{notation}

\newaliascnt{construction}{definition}
\newtheorem{construction}[construction]{Construction}
\aliascntresetthe{construction}

% Example: document-wide numbering (not per section / chapter)
\newtheorem{example}{Example}

\theoremstyle{remark}
\newtheorem{remark}{Remark}[section]
\newtheorem*{remark*}{Remark}

\newaliascnt{heuristic}{remark}
\newtheorem{heuristic}[heuristic]{Heuristic}
\aliascntresetthe{heuristic}

\newaliascnt{convention}{remark}
\newtheorem{convention}[convention]{Convention}
\aliascntresetthe{convention}

\newaliascnt{warning}{remark}
\newtheorem{warning}[warning]{Warning}
\aliascntresetthe{warning}

\newaliascnt{proofsketch}{remark}
\newtheorem{proofsketch}[proofsketch]{Proof sketch}
\aliascntresetthe{proofsketch}

\newaliascnt{checkpoint}{remark}
\newtheorem{checkpoint}[checkpoint]{Checkpoint}
\aliascntresetthe{checkpoint}

% Note: document-wide numbering (not per section)
\newtheorem{note}{Note}
\newtheorem*{note*}{Note}

\newaliascnt{examnote}{note}
\newtheorem{examnote}[examnote]{Exam note}
\aliascntresetthe{examnote}

\newaliascnt{study}{note}
\newtheorem{study}[study]{Study note}
\aliascntresetthe{study}

\newaliascnt{summary}{note}
\newtheorem{summary}[summary]{Summary}
\aliascntresetthe{summary}

\newtheorem{examproblem}{Problem}[section]
\newtheorem*{examproblem*}{Problem}

% Bilingual aliases
\newenvironment{dingli}{\begin{theorem}}{\end{theorem}}
\newenvironment{yinli}{\begin{lemma}}{\end{lemma}}
\newenvironment{mingti}{\begin{proposition}}{\end{proposition}}
\newenvironment{tuilun}{\begin{corollary}}{\end{corollary}}
\newenvironment{dingyi}{\begin{definition}}{\end{definition}}
\newenvironment{li}{\begin{example}}{\end{example}}
\newenvironment{zhu}{\begin{remark}}{\end{remark}}
\newenvironment{biji}{\begin{note}}{\end{note}}
\newenvironment{kaoshi}{\begin{examnote}}{\end{examnote}}

\renewcommand{\qedsymbol}{\ensuremath{\square}}

% PDF sidebar outline + printed TOC for QE exam transcripts.
% Load in the preamble AFTER hyperref and AFTER \newcommand{\examstatement}.
\hypersetup{
  unicode=true,
  bookmarks=true,
  bookmarksopen=true,
  bookmarksopenlevel=3,
  bookmarksnumbered=false,
  bookmarksdepth=3
}
\IfFileExists{bookmark.sty}{%
  \usepackage{bookmark}%
  \bookmarksetup{open,openlevel=3,numbered=false,depth=3}%
}{}
% Unnumbered in print, but still written to the TOC / PDF outline
% down to \subsubsection. Starred headings create neither.
\setcounter{secnumdepth}{-1}
\setcounter{tocdepth}{3}
\makeatletter
\@ifundefined{examstatement}{}{%
  \let\qe@origexamstatement\examstatement
  \renewcommand{\examstatement}[1]{%
    \par\phantomsection
    \addcontentsline{toc}{subsection}{#1}%
    \qe@origexamstatement{#1}%
  }%
}
\makeatother


\title{\textbf{QUALIFYING EXAM: APPLIED MATHEMATICS}\\[0.5em]\Large FALL 2024}
\date{}
\author{}

\begin{document}

\maketitle
\vspace{-3em}

\tableofcontents

\section{Statement of the problems}

\examstatement{Problem 1. [20 points]}
The \(n\)-th Chebyshev polynomial is defined by: \(T_{n}(x) = \cos(n \cdot \arccos(x))\). Answer the following questions related to \(T_{n}(x)\).
\begin{enumerate}[label=(\roman*)]
    \item (6 points) Prove that \(T_{n}(x)\) is a polynomial of degree \(n\).
    \item (6 points) For any \(x\) with \(0 < x < 1\), show that \(T_{n}(2x - 1) = T_{2n}(\sqrt{x})\).
    \item (8 points) Prove that the system of Chebyshev polynomials \(\{T_{k} : 0 \le k < n\}\) is orthogonal with respect to the discrete inner product \((u, v) = \sum_{k=1}^{n} u(x_{k}) v(x_{k})\), where \(\{x_{k}\}\) are the Chebyshev points \(x_{k} = \cos\bigl(\frac{2k-1}{2n}\pi\bigr)\).
\end{enumerate}

\examstatement{Problem 2. [15 points]}
Let \(A\) be a symmetric positive definite matrix. Assume that the conjuage gradient method is applied to solve the linear system \(Ax = b\), where \(x^{*}\) is the exact solution.
\begin{enumerate}[label=(\roman*)]
    \item (10 points) Prove the following error estimate:
    \[
    \|x_{k} - x^{*}\|_{A} \le \left( \frac{\sqrt{\kappa_{2}(A)} - 1}{\sqrt{\kappa_{2}(A)} + 1} \right)^{k} \|x_{0} - x^{*}\|_{A},
    \]
    where \(x_{k}\) is the solution obtained by the \(k\)-th iteration, \(\|\cdot\|_{A}\) denotes the \(A\)-norm defined by \(\|x\|_{A} = \sqrt{x^{*} A x}\), and \(\kappa_{2}(A) = \dfrac{\max \lambda(A)}{\min \lambda(A)}\) denotes the condition number of \(A\) under the \(l_{2}\) norm. (If you apply any theorem from approximation theory, please state the theorem clearly. No need to prove that theorem.)
    \item (5 points) Describe one Krylov subspace method for non-Hermitian matrices (\(A \neq A^{*}\)). (You can use pseudo-code or words/formulas as long as the steps are clear.)
\end{enumerate}

\examstatement{Problem 3. [15 points]}
Consider Kepler's equation

\[
f(x) = x - \epsilon \sin x - \eta, \quad (0 < |\epsilon| < 1),\quad \eta \in \mathbb{R}.
\]

\begin{enumerate}[label=(\roman*)]
    \item (5 points) Show that for each \(\epsilon, \eta\), there is exactly one real root \(\alpha = \alpha(\epsilon, \eta)\). Furthermore, \(\eta - |\epsilon| \le \alpha \le \eta + |\epsilon|\).
    \item (5 points) Writing the equation in fixed point form:
    \[
    x = \varphi(x), \qquad \varphi(x) = \epsilon \sin x + \eta,
    \]
    show that the fixed point iteration \(x_{n+1} = \varphi(x_{n})\) converges for arbitrary staring value \(x_{0}\).
    \item (5 points) Let \(m\) be an integer such that \(m\pi < \eta < (m+1)\pi\). Show that Newton's method with starting value
    \[
    x_{0} =
    \begin{cases}
    (m+1)\pi & \text{if } (-1)^{m} \epsilon > 0, \\
    m\pi & \text{otherwise,}
    \end{cases}
    \]
    is guaranteed to converge (monotonically) to \(\alpha(\epsilon, \eta)\).
\end{enumerate}

\examstatement{Problem 4. [15 points]}
For the system

\[
u_{t} = v_{x}, \qquad v_{t} = u_{x},
\]

analyze the truncation error and stability of the scheme

\[
\frac{1}{\tau} \left( u_{j}^{n+1} - \frac12 (u_{j+1}^{n} + u_{j-1}^{n}) \right) = \frac{1}{2h} (v_{j+1}^{n} - v_{j-1}^{n}),
\]

\[
\frac{1}{\tau} \left( v_{j}^{n+1} - \frac12 (v_{j+1}^{n} + v_{j-1}^{n}) \right) = \frac{1}{2h} (u_{j+1}^{n} - u_{j-1}^{n}).
\]

\examstatement{Problem 5. [15 points]}
Write and prove the maximum principle of the centered finite difference scheme for discretizing the equation

\[
u_{xx} + u_{yy} + d(x,y) u_{x} + e(x,y) u_{y} + f(x,y) u = 0, \quad f < 0,
\]

under some suitable assumptions.

\noindent\textbf{Remark: You can choose either (6) or (7). The points will be decided as \(\max((6), (7))\).}

\examstatement{Problem 6. [20 points]}
Consider the following boundary value problem for \(y = y(x)\) on \([0, 1]\) as \(0 < \epsilon \ll 1\),

\[
\epsilon y'' + \epsilon(1+x)^{2} y' - y = x - 1, \qquad y(0) = \alpha, \quad y(1) = -1.
\]

\begin{enumerate}[label=(\roman*)]
    \item (10 points) Suppose \(\alpha = 1\). Construct a composite expansion of the above problem and sketch the solution.
    \item (5 points) Construct a composite expansion of the above problem for \(\alpha = 0\).
    \item (5 points) What is the accuracy of your solution in \(\epsilon\)? Formally explain your conclusion. Consider the first case only.
\end{enumerate}

\examstatement{Problem 7. [20 points]}
Let \(f : \mathbb{R}^{n} \to \mathbb{R}\) be convex and define \(\partial f(x)\) to be the subgradient set of \(f\) at \(x\).
\begin{enumerate}[label=(\roman*)]
    \item (10 points) If \(x \in \operatorname{int} \operatorname{dom} f\), prove that \(\partial f(x)\) is nonempty and bounded.
    \item (5 points) Assume \(f = \|x\|_{1} = \sum_{i=1}^{n} |x_{i}|\), write down the \(\partial f(x)\).
    \item (5 points) Given \(y \in \mathbb{R}^{n}\) and \(\lambda > 0\), calculate the closed form solution of the problem:
    \[
    \min_{x} \tfrac12 \|x - y\|_{2}^{2} + \lambda \|x\|_{1}.
    \]
\end{enumerate}

\noindent\rule{\linewidth}{0.4pt}

\section{Statement and solutions}

\subsection{Problem 1. [20 points]}
The \(n\)-th Chebyshev polynomial is defined by: \(T_{n}(x) = \cos(n \cdot \arccos(x))\). Answer the following questions related to \(T_{n}(x)\).
\begin{enumerate}[label=(\roman*)]
    \item (6 points) Prove that \(T_{n}(x)\) is a polynomial of degree \(n\).
    \item (6 points) For any \(x\) with \(0 < x < 1\), show that \(T_{n}(2x - 1) = T_{2n}(\sqrt{x})\).
    \item (8 points) Prove that the system of Chebyshev polynomials \(\{T_{k} : 0 \le k < n\}\) is orthogonal with respect to the discrete inner product \((u, v) = \sum_{k=1}^{n} u(x_{k}) v(x_{k})\), where \(\{x_{k}\}\) are the Chebyshev points \(x_{k} = \cos\bigl(\frac{2k-1}{2n}\pi\bigr)\).
\end{enumerate}

\subsection{Problem 2. [15 points]}
Let \(A\) be a symmetric positive definite matrix. Assume that the conjuage gradient method is applied to solve the linear system \(Ax = b\), where \(x^{*}\) is the exact solution.
\begin{enumerate}[label=(\roman*)]
    \item (10 points) Prove the following error estimate:
    \[
    \|x_{k} - x^{*}\|_{A} \le \left( \frac{\sqrt{\kappa_{2}(A)} - 1}{\sqrt{\kappa_{2}(A)} + 1} \right)^{k} \|x_{0} - x^{*}\|_{A},
    \]
    where \(x_{k}\) is the solution obtained by the \(k\)-th iteration, \(\|\cdot\|_{A}\) denotes the \(A\)-norm defined by \(\|x\|_{A} = \sqrt{x^{*} A x}\), and \(\kappa_{2}(A) = \dfrac{\max \lambda(A)}{\min \lambda(A)}\) denotes the condition number of \(A\) under the \(l_{2}\) norm. (If you apply any theorem from approximation theory, please state the theorem clearly. No need to prove that theorem.)
    \item (5 points) Describe one Krylov subspace method for non-Hermitian matrices (\(A \neq A^{*}\)). (You can use pseudo-code or words/formulas as long as the steps are clear.)
\end{enumerate}

\subsection{Problem 3. [15 points]}
Consider Kepler's equation

\[
f(x) = x - \epsilon \sin x - \eta, \quad (0 < |\epsilon| < 1),\quad \eta \in \mathbb{R}.
\]

\begin{enumerate}[label=(\roman*)]
    \item (5 points) Show that for each \(\epsilon, \eta\), there is exactly one real root \(\alpha = \alpha(\epsilon, \eta)\). Furthermore, \(\eta - |\epsilon| \le \alpha \le \eta + |\epsilon|\).
    \item (5 points) Writing the equation in fixed point form:
    \[
    x = \varphi(x), \qquad \varphi(x) = \epsilon \sin x + \eta,
    \]
    show that the fixed point iteration \(x_{n+1} = \varphi(x_{n})\) converges for arbitrary staring value \(x_{0}\).
    \item (5 points) Let \(m\) be an integer such that \(m\pi < \eta < (m+1)\pi\). Show that Newton's method with starting value
    \[
    x_{0} =
    \begin{cases}
    (m+1)\pi & \text{if } (-1)^{m} \epsilon > 0, \\
    m\pi & \text{otherwise,}
    \end{cases}
    \]
    is guaranteed to converge (monotonically) to \(\alpha(\epsilon, \eta)\).
\end{enumerate}

\subsection{Problem 4. [15 points]}
For the system

\[
u_{t} = v_{x}, \qquad v_{t} = u_{x},
\]

analyze the truncation error and stability of the scheme

\[
\frac{1}{\tau} \left( u_{j}^{n+1} - \frac12 (u_{j+1}^{n} + u_{j-1}^{n}) \right) = \frac{1}{2h} (v_{j+1}^{n} - v_{j-1}^{n}),
\]

\[
\frac{1}{\tau} \left( v_{j}^{n+1} - \frac12 (v_{j+1}^{n} + v_{j-1}^{n}) \right) = \frac{1}{2h} (u_{j+1}^{n} - u_{j-1}^{n}).
\]

\subsection{Problem 5. [15 points]}
Write and prove the maximum principle of the centered finite difference scheme for discretizing the equation

\[
u_{xx} + u_{yy} + d(x,y) u_{x} + e(x,y) u_{y} + f(x,y) u = 0, \quad f < 0,
\]

under some suitable assumptions.

\noindent\textbf{Remark: You can choose either (6) or (7). The points will be decided as \(\max((6), (7))\).}

\subsection{Problem 6. [20 points]}
Consider the following boundary value problem for \(y = y(x)\) on \([0, 1]\) as \(0 < \epsilon \ll 1\),

\[
\epsilon y'' + \epsilon(1+x)^{2} y' - y = x - 1, \qquad y(0) = \alpha, \quad y(1) = -1.
\]

\begin{enumerate}[label=(\roman*)]
    \item (10 points) Suppose \(\alpha = 1\). Construct a composite expansion of the above problem and sketch the solution.
    \item (5 points) Construct a composite expansion of the above problem for \(\alpha = 0\).
    \item (5 points) What is the accuracy of your solution in \(\epsilon\)? Formally explain your conclusion. Consider the first case only.
\end{enumerate}

\subsection{Problem 7. [20 points]}
Let \(f : \mathbb{R}^{n} \to \mathbb{R}\) be convex and define \(\partial f(x)\) to be the subgradient set of \(f\) at \(x\).
\begin{enumerate}[label=(\roman*)]
    \item (10 points) If \(x \in \operatorname{int} \operatorname{dom} f\), prove that \(\partial f(x)\) is nonempty and bounded.
    \item (5 points) Assume \(f = \|x\|_{1} = \sum_{i=1}^{n} |x_{i}|\), write down the \(\partial f(x)\).
    \item (5 points) Given \(y \in \mathbb{R}^{n}\) and \(\lambda > 0\), calculate the closed form solution of the problem:
    \[
    \min_{x} \tfrac12 \|x - y\|_{2}^{2} + \lambda \|x\|_{1}.
    \]
\end{enumerate}

\end{document}
